% =============================================================================
% Part 4 — fragment for the "Towards Adaptive and Situated AI..." paper version.
% Narrative register to match the System Evaluation section (plain language,
% rounded figures, no \texttt, ties back to the expert interviews). Numbers
% verified against eval/results/part4_20260618_115438.{json,summary.txt}.
%
% ONE edit: REPLACE the whole "\subsection{TODO: System Personalization
% Evaluation}" block (its header and the single existing paragraph) with the
% block below. The first paragraph is kept verbatim from the existing draft.
% =============================================================================

\subsection{System Personalization Evaluation}

This evaluation phase directly addresses one of the key concerns identified in the expert interviews: an AI guide is only acceptable in a cultural institution if it does not fabricate historical facts and if its responses remain accurate across different personalization settings, always grounded in the museum's own data.

Because GuiA personalizes by prompting, the explanation style, the visitor profile, the worked examples, and even the order in which the retrieved facts and the instructions appear are all prompt choices. Prior work warns that such choices can change a model's output in ways that are easy to overlook, and recommends comparing prompt variants on a fixed set of questions rather than trusting a single manual trial \cite{PromptSurvey2024}. We therefore measured how much each choice actually changes the answer, and whether that change is a genuine effect or only the random variation the model shows when asked the same thing twice.

The test holds everything constant except the prompt. For each question the museum information is retrieved once and reused for every version of the prompt, so that any difference in the answer can only come from the prompt itself and not from the retrieval step. A baseline prompt, identical to the one used by the deployed system, is compared against three variants that each change a single element: the worked examples are removed, the explanation style is switched from the standard guide to the more scholarly one, and the block of retrieved facts is moved from after the instructions to before them. Each version is run several times, and the difference between a variant and the baseline is read against the variation observed between repeated runs of the same prompt. A value near that baseline variation means the change made no real difference; a clearly larger value means it did. The two kinds of change are read in opposite directions: switching the explanation style is meant to change the answer, so a large difference is desirable, whereas removing the examples or moving the block of facts should not change the substance of the answer, so here a small difference is what we want. Twelve grounded questions were used in each language, repeated across the baseline and the three variants.

The clearest result is that the worked examples matter. Removing them changed the answer about twice as much as the model's own run-to-run variation, and the direction was consistent: answers written with the examples were longer and more explanatory, while answers written without them were noticeably terser, often a single sentence. This confirms, on actual museum questions, the rationale for including them. They do not only teach the guide when to decline; they also shape how much it explains. Switching the explanation style, by contrast, changed the answer only modestly, not much more than the natural variation between repeated runs. The style setting therefore has a real but limited effect on the generated text and is a candidate for strengthening if a sharper contrast between styles is desired. Moving the block of retrieved facts produced a similarly small change, which is the desired behaviour for an element that should not alter the substance of an answer. Importantly, because the retrieved facts were held fixed, none of the personalization settings caused the guide to drift away from the museum data: personalization changed the wording and the level of detail, not the underlying facts, which is exactly the property the museum requires.

These results should be read as a first indication. The measure compares answers by the words they share, so it detects that an answer changed but not whether its meaning was preserved, and a faithful rephrasing counts as a difference; it is therefore a test of stability and sensitivity rather than of correctness. A small number of empty or malformed generations were set aside, and the sample is modest, twelve questions per language, so the figures point to clear tendencies rather than precise values. Even so, the test supports the personalization design: the elements meant to keep the guide reliable have a visible effect, the incidental choices do not disturb the facts, and the style setting works but could be made stronger.
